% =============================================================================
% APPENDIX FEH: LIT-FEHR: Ernst Fehr Research Integration
% =============================================================================
% Module: BCM2_LIT_FEHR
% Version: 1.0 (January 2026)
% Status: Final
% Language: English
%
% Comprehensive documentation of Ernst Fehr's scientific contributions,
% their integration into the Evidence-Based Framework (EBF), extractable
% parameters, cases, and practical implications for practitioners.
%
% SSOT Reference: data/researcher-registry.yaml (RES-FEHR-E)
% =============================================================================

% =============================================================================
% CROSS-REFERENCE STRUCTURE
% =============================================================================
%
% CHAPTER LINKAGE:
% - Primary Chapter: Chapter 10 (CORE-WHAT: Social Preferences)
% - Secondary Chapters: Chapter 5 (CORE-WHO), Chapter 11 (CORE-HOW), Chapter 9 (CORE-WHEN)
%
% DEPENDENCIES:
% - Appendix C (CORE-WHAT): Utility dimensions U^S
% - Appendix BBB (CORE-WHERE): Parameter repository
% - Appendix B (CORE-HOW): Complementarity γ
%
% DEPENDENTS:
% - Appendix FRM (LIT-FEHR-METHOD): Methodological integration
% - All DOMAIN appendices using social preference parameters
%
% =============================================================================

\begin{tcolorbox}[colback=purple!5!white, colframe=purple!75!black,
    title=Chapter Linkage]

\textbf{Primary Chapter:} Chapter 10: Social Preferences and Other-Regarding Behavior
\begin{itemize}[nosep]
\item Section 10.2: Inequity Aversion $\leftrightarrow$ This Appendix Section \ref{sec:feh-inequity}
\item Section 10.3: Strong Reciprocity $\leftrightarrow$ This Appendix Section \ref{sec:feh-reciprocity}
\item Section 10.4: Gift Exchange $\leftrightarrow$ This Appendix Section \ref{sec:feh-gift}
\end{itemize}

\textbf{Secondary Chapters:}
\begin{itemize}[nosep]
\item Chapter 5: CORE-WHO --- preference type heterogeneity
\item Chapter 9: CORE-WHEN --- context effects on social preferences
\item Chapter 11: CORE-HOW --- complementarity in cooperation mechanisms
\end{itemize}

\textbf{Reading Guide:}
\begin{itemize}[nosep]
\item For conceptual overview: Read Chapter 10 first
\item For formal parameter values: This appendix provides calibration
\item For methodology: See Appendix FRM (LIT-FEHR-METHOD)
\end{itemize}

\end{tcolorbox}


\chapter{LIT-FEHR: Ernst Fehr Research Integration}
\label{app:feh}


\begin{tcolorbox}[colback=blue!5!white, colframe=blue!75!black,
    title=Quick Reference: Appendix FEH]

\textbf{Appendix:} FEH (LIT-FEHR)

\textbf{Registry ID:} RES-FEHR-E

\textbf{Core Question:} How do Ernst Fehr's contributions enable behavioral economics
in EBF, and which parameters can be estimated from his research?

\textbf{Key Concepts:}
\begin{itemize}[nosep]
\item Inequity Aversion ($\alpha$, $\beta$)
\item Strong Reciprocity
\item Gift Exchange
\item Preference Heterogeneity
\end{itemize}

\textbf{Key Parameters:}
\begin{align}
\alpha &= 0.85 \quad [0.5, 4.0] \quad \text{(disadvantageous inequity aversion)} \\
\beta &= 0.315 \quad [0.0, 0.6] \quad \text{(advantageous inequity aversion)} \\
\gamma_{\text{trust,contract}} &> 0 \quad \text{(complementarity in cooperation)}
\end{align}

\textbf{Cross-References:} C (CORE-WHAT), BBB (CORE-WHERE), FRM (METHOD), B (CORE-HOW)
\end{tcolorbox}


\begin{abstract}
This appendix documents the comprehensive research contributions of Ernst Fehr
(h-index: 157, 180,000+ citations) to behavioral economics and their integration
into the Evidence-Based Framework (EBF). Fehr's work provides the empirical
foundation for EBF's treatment of social preferences, establishing that utility
includes other-regarding components ($U^S$) with measurable parameters. The appendix
catalogs 144 integrated papers, documents non-integrated papers with evaluation,
and provides practitioners with actionable parameter estimates and case examples.
\end{abstract}


\begin{tcolorbox}[
    colback=orange!5!white,
    colframe=orange!75!black,
    title={\textbf{Appendix Scope: Ziel / In-Scope / Out-of-Scope / Constraints}},
    fonttitle=\bfseries
]

\textbf{Ziel (Objective):}
\begin{itemize}[nosep]
    \item Document Ernst Fehr's scientific contributions and their integration into EBF
\end{itemize}

\textbf{In-Scope:}
\begin{itemize}[nosep]
    \item Scientific contributions to economics (5 key areas)
    \item EBF framework contributions (5 CORE dimensions)
    \item Extractable parameters with confidence intervals
    \item Cases and worked examples from Fehr's papers
    \item Practical implications for behavior change practitioners
    \item Paper bibliography (integrated and non-integrated with evaluation)
\end{itemize}

\textbf{Out-of-Scope (delegiert an andere Appendices):}
\begin{itemize}[nosep]
    \item Formal derivations of Fehr-Schmidt model $\rightarrow$ Appendix D (FORMAL)
    \item Methodological critique (integration vs. falsifiability) $\rightarrow$ Appendix FRM
    \item Parameter repository details $\rightarrow$ Appendix BBB (CORE-WHERE)
    \item Context framework details $\rightarrow$ Appendix V (CORE-WHEN)
\end{itemize}

\textbf{Constraints (Anwendungsgrenzen):}
\begin{itemize}[nosep]
    \item Parameters are population averages; individual variation is substantial
    \item Laboratory estimates may differ from field contexts (external validity)
    \item Neuroeconomics papers excluded from EBF scope
\end{itemize}

\textbf{Lieferobjekte (Deliverables):}
\begin{itemize}[nosep]
    \item Table of extractable parameters with sources
    \item Type distribution for population heterogeneity
    \item Worked examples for practitioner application
    \item Complete paper bibliography with integration decisions
\end{itemize}

\end{tcolorbox}


%%%%%%%%%%%%%%%%%%%%%%%%%%%%%%%%%%%%%%%%%%%%%%%%%%%%%%%%%%%%%%%%%%%%%%%%%%%%%%%
\section{The Fundamental Question}
\label{sec:feh-fundamental}
%%%%%%%%%%%%%%%%%%%%%%%%%%%%%%%%%%%%%%%%%%%%%%%%%%%%%%%%%%%%%%%%%%%%%%%%%%%%%%%

\subsection{Why This Matters}

Ernst Fehr is among the most influential economists of the past three decades,
with an h-index of 157 and over 180,000 citations. His research has fundamentally
challenged the homo economicus assumption that dominated economics, demonstrating
empirically that:

\begin{enumerate}
    \item People care about fairness, not just their own payoffs
    \item Cooperation can emerge and persist even among strangers
    \item Context (institutions, norms, culture) shapes behavior systematically
    \item Population heterogeneity in preferences explains seemingly contradictory findings
\end{enumerate}

For the Evidence-Based Framework, Fehr's contributions are foundational. Without
his empirical parameters ($\alpha$, $\beta$), EBF cannot calibrate social preference
models. Without his experimental methods, EBF cannot validate predictions.

\subsection{The Core Problem}

How do we systematically integrate 343 publications spanning 35+ years into a
coherent framework for behavioral interventions?

\begin{tcolorbox}[colback=red!5!white, colframe=red!75!black,
    title=The Fundamental Question]
\textbf{Which of Fehr's contributions provide extractable parameters, testable
theories, or replicable cases for EBF---and which remain outside scope?}
\end{tcolorbox}


%%%%%%%%%%%%%%%%%%%%%%%%%%%%%%%%%%%%%%%%%%%%%%%%%%%%%%%%%%%%%%%%%%%%%%%%%%%%%%%
\section{Research Principles (Axioms)}
\label{sec:feh-axioms}
%%%%%%%%%%%%%%%%%%%%%%%%%%%%%%%%%%%%%%%%%%%%%%%%%%%%%%%%%%%%%%%%%%%%%%%%%%%%%%%

Fehr's research program is built on foundational methodological principles that
guide EBF integration:

\begin{itemize}
    \item \textbf{AX-F1 (Experimental Validity):} Behavioral claims require controlled
    experimental evidence, not just observational data or introspection.

    \item \textbf{AX-F2 (Incentive Compatibility):} Experimental subjects must face
    real monetary consequences; hypothetical choices systematically differ from
    consequential ones.

    \item \textbf{AX-F3 (Heterogeneity):} Population distributions of preferences
    (e.g., $\alpha$, $\beta$) must be estimated, not assumed homogeneous.

    \item \textbf{AX-F4 (Context Sensitivity):} Social preferences vary with
    institutional context ($\Psi$); lab findings require field validation.

    \item \textbf{AX-F5 (Parsimony):} Models should use minimal parameters to explain
    observed behavior; the Fehr-Schmidt model uses only $\alpha$ and $\beta$.
\end{itemize}


%%%%%%%%%%%%%%%%%%%%%%%%%%%%%%%%%%%%%%%%%%%%%%%%%%%%%%%%%%%%%%%%%%%%%%%%%%%%%%%
\section{Key Results}
\label{sec:feh-results}
%%%%%%%%%%%%%%%%%%%%%%%%%%%%%%%%%%%%%%%%%%%%%%%%%%%%%%%%%%%%%%%%%%%%%%%%%%%%%%%

The following empirically validated results form the foundation for EBF parameter
estimation:

\begin{enumerate}
    \item \textbf{R1:} Approximately 40\% of subjects exhibit inequity aversion
    ($\alpha > 0$ or $\beta > 0$) in standard dictator and ultimatum games.

    \item \textbf{R2:} The median value of $\alpha = 0.85$ (disadvantageous inequity
    aversion) is robust across 15+ countries and 50+ replications.

    \item \textbf{R3:} Strong reciprocity (costly punishment of defectors) persists
    even in one-shot anonymous interactions, contradicting pure self-interest.

    \item \textbf{R4:} Gift exchange effects on effort provision range from 10-40\%
    above equilibrium predictions in lab and field settings.

    \item \textbf{R5:} Context ($\Psi$) modulates social preferences: competitive
    framing reduces cooperation by 20-30\% compared to cooperative framing.
\end{enumerate}


%%%%%%%%%%%%%%%%%%%%%%%%%%%%%%%%%%%%%%%%%%%%%%%%%%%%%%%%%%%%%%%%%%%%%%%%%%%%%%%
\section{Scientific Contributions to Economics}
\label{sec:feh-economics}
%%%%%%%%%%%%%%%%%%%%%%%%%%%%%%%%%%%%%%%%%%%%%%%%%%%%%%%%%%%%%%%%%%%%%%%%%%%%%%%

Ernst Fehr's contributions to economics span five major areas, each with
lasting impact on theory and practice.

%------------------------------------------------------------------------------
\subsection{Inequity Aversion Model (Fehr-Schmidt 1999)}
\label{sec:feh-inequity}
%------------------------------------------------------------------------------

The Fehr-Schmidt model \citep{fehr1999theory} fundamentally redefined utility
by incorporating distributional concerns:

\begin{equation}
U_i(x) = x_i - \alpha_i \cdot \frac{1}{n-1} \sum_{j \neq i} \max\{x_j - x_i, 0\}
       - \beta_i \cdot \frac{1}{n-1} \sum_{j \neq i} \max\{x_i - x_j, 0\}
\label{eq:fehr-schmidt}
\end{equation}

Where:
\begin{itemize}[nosep]
    \item $x_i$ = material payoff to agent $i$
    \item $\alpha_i$ = sensitivity to \textit{disadvantageous} inequity (being behind)
    \item $\beta_i$ = sensitivity to \textit{advantageous} inequity (being ahead)
\end{itemize}

\textbf{Key constraint:} $\alpha_i \geq \beta_i$ and $0 \leq \beta_i < 1$

\paragraph{Impact.} With 19,000+ citations, this is the most cited behavioral
economics model. It explains:
\begin{itemize}[nosep]
    \item Ultimatum game rejections (even with real stakes)
    \item Public goods provision with punishment
    \item Wage compression in labor markets
    \item Resistance to ``unfair'' price increases
\end{itemize}

%------------------------------------------------------------------------------
\subsection{Strong Reciprocity}
\label{sec:feh-reciprocity}
%------------------------------------------------------------------------------

Fehr's work on strong reciprocity \citep{fehr2002strong, fehr2002altruistic}
demonstrates that humans cooperate and punish defectors even when:
\begin{itemize}[nosep]
    \item The interaction is one-shot (no reputation effects)
    \item Partners are anonymous (no social pressure)
    \item Punishment is costly to the punisher
\end{itemize}

\textbf{Definition:} Strong reciprocity = propensity to cooperate with cooperators
and punish defectors, even at personal cost, in anonymous one-shot interactions.

\paragraph{Key finding.} In public goods games with punishment, cooperation
stabilizes at 80-95\%, compared to 20-40\% without punishment. The presence
of strong reciprocators transforms group dynamics.

%------------------------------------------------------------------------------
\subsection{Gift Exchange and Efficiency Wages}
\label{sec:feh-gift}
%------------------------------------------------------------------------------

The gift exchange literature \citep{fehr1993fairness, fehr1998reciprocity}
provides experimental validation for efficiency wage theory:

\begin{itemize}[nosep]
    \item Workers respond to higher wages with higher effort (reciprocity)
    \item This effect persists even without monitoring or reputation
    \item Firms can sustain above-market wages profitably
\end{itemize}

\textbf{EBF implication:} Monetary incentives interact with social preferences.
A wage increase can be a ``gift'' eliciting reciprocal effort, not just a
price for labor.

%------------------------------------------------------------------------------
\subsection{Experimental Methods}
\label{sec:feh-methods}
%------------------------------------------------------------------------------

Fehr pioneered and standardized key experimental paradigms:

\begin{table}[htbp]
\centering
\caption{Experimental Games Standardized by Fehr's Research}
\label{tab:feh-games}
\begin{tabular}{@{}lll@{}}
\toprule
\textbf{Game} & \textbf{Measures} & \textbf{Key Finding} \\
\midrule
Ultimatum Game & Fairness rejection threshold & Modal offer 40-50\% \\
Dictator Game & Pure altruism/fairness & Mean transfer 20-30\% \\
Trust Game & Trust and trustworthiness & 50\% trusted, 30\% returned \\
Public Goods + Punishment & Cooperation sustainability & 80-95\% with punishment \\
Gift Exchange Game & Wage-effort reciprocity & Positive correlation \\
\bottomrule
\end{tabular}
\end{table}

%------------------------------------------------------------------------------
\subsection{Behavioral Contract Theory}
\label{sec:feh-contracts}
%------------------------------------------------------------------------------

More recent work \citep{fehr2007competition} integrates behavioral assumptions
into mechanism design, showing that optimal contracts must account for:
\begin{itemize}[nosep]
    \item Heterogeneity in social preferences
    \item Reference-dependent utility
    \item Loss aversion in effort provision
\end{itemize}


%%%%%%%%%%%%%%%%%%%%%%%%%%%%%%%%%%%%%%%%%%%%%%%%%%%%%%%%%%%%%%%%%%%%%%%%%%%%%%%
\section{Contributions to the EBF Framework}
\label{sec:feh-ebf}
%%%%%%%%%%%%%%%%%%%%%%%%%%%%%%%%%%%%%%%%%%%%%%%%%%%%%%%%%%%%%%%%%%%%%%%%%%%%%%%

Fehr's research maps onto multiple CORE dimensions of EBF:

%------------------------------------------------------------------------------
\subsection{CORE-WHAT: Social Utility Component}
%------------------------------------------------------------------------------

Fehr establishes that utility is not unidimensional:

\begin{equation}
U = U^F + U^S + U^P + U^E + U^{ID}
\end{equation}

His work provides the foundation for $U^S$ (social utility), demonstrating:
\begin{itemize}[nosep]
    \item People have preferences over distributions, not just own payoffs
    \item Fairness concerns are measurable via $\alpha$ and $\beta$
    \item Social comparison is a fundamental driver of satisfaction
\end{itemize}

\textbf{EBF Integration:} Appendix C (CORE-WHAT) incorporates Fehr-Schmidt
parameters as default social preference calibration.

%------------------------------------------------------------------------------
\subsection{CORE-HOW: Complementarity in Cooperation}
%------------------------------------------------------------------------------

Fehr's work reveals positive complementarity ($\gamma > 0$) between cooperation
mechanisms:

\begin{equation}
\gamma(\text{trust}, \text{contract}) > 0
\end{equation}

Trust and formal contracts are \textit{complements}, not substitutes. The
recent Nature paper on ``Super-additive Cooperation'' \citep{efferson2024superadditive}
demonstrates:

\begin{equation}
\gamma(\text{repetition}, \text{competition}) > 0 \quad \text{(super-additive)}
\end{equation}

\textbf{EBF Integration:} Appendix B (CORE-HOW) uses these findings to justify
bundled interventions.

%------------------------------------------------------------------------------
\subsection{CORE-WHEN: Context Effects}
%------------------------------------------------------------------------------

Fehr's banking culture research \citep{cohn2014business, cohn2015evidence}
demonstrates systematic context effects:

\begin{itemize}[nosep]
    \item $\Psi_{\text{norm}}$: Social norms shift behavior significantly
    \item $\Psi_{\text{culture}}$: Professional identity affects honesty
    \item $\Psi_{\text{inst}}$: Institutional rules shape cooperation
\end{itemize}

\textbf{Key finding:} The \textit{same} individuals behave differently when
their professional (vs. personal) identity is salient.

%------------------------------------------------------------------------------
\subsection{CORE-WHO: Preference Heterogeneity}
%------------------------------------------------------------------------------

Fehr's work reveals that the population is not a representative agent but
a \textit{mixture of types}:

\begin{table}[htbp]
\centering
\caption{Type Distribution from Fehr's Research}
\label{tab:feh-types}
\begin{tabular}{@{}lccc@{}}
\toprule
\textbf{Type} & \textbf{Proportion} & \textbf{$\alpha$} & \textbf{$\beta$} \\
\midrule
Selfish & 30\% & 0 & 0 \\
Moderate & 30\% & 0.5 & 0.25 \\
Fair & 30\% & 1.0 & 0.5 \\
Strongly Fair & 10\% & $>1$ & $>0.5$ \\
\bottomrule
\end{tabular}
\end{table}

\textbf{EBF Implication:} Interventions must account for heterogeneity.
A policy optimal for ``fair'' types may backfire for ``selfish'' types.


%%%%%%%%%%%%%%%%%%%%%%%%%%%%%%%%%%%%%%%%%%%%%%%%%%%%%%%%%%%%%%%%%%%%%%%%%%%%%%%
\section{Extractable Parameters}
\label{sec:feh-parameters}
%%%%%%%%%%%%%%%%%%%%%%%%%%%%%%%%%%%%%%%%%%%%%%%%%%%%%%%%%%%%%%%%%%%%%%%%%%%%%%%

The following parameters can be reliably estimated from Fehr's research:

\begin{table}[htbp]
\centering
\caption{EBF Parameters from Fehr's Research}
\label{tab:feh-params}
\renewcommand{\arraystretch}{1.4}
\begin{tabular}{@{}lccll@{}}
\toprule
\textbf{Parameter} & \textbf{Point Est.} & \textbf{Range} & \textbf{Source} & \textbf{Appendix} \\
\midrule
$\alpha$ (disadvantageous) & 0.85 & [0.5, 4.0] & \citet{fehr1999theory} & BBB \\
$\beta$ (advantageous) & 0.315 & [0.0, 0.6] & \citet{fehr1999theory} & BBB \\
Punishment cost ratio & 3:1 & [2:1, 4:1] & \citet{fehr2002strong} & BBB \\
Trust rate (baseline) & 0.50 & [0.40, 0.60] & Trust Game meta & BBB \\
Reciprocity rate & 0.30 & [0.20, 0.40] & Trust Game meta & BBB \\
Gift exchange elasticity & 0.25 & [0.15, 0.35] & \citet{fehr1998reciprocity} & BBB \\
\bottomrule
\end{tabular}
\end{table}

\textbf{Usage note:} These are population averages. For segment-specific
calibration, use the type distribution in Table \ref{tab:feh-types}.


%%%%%%%%%%%%%%%%%%%%%%%%%%%%%%%%%%%%%%%%%%%%%%%%%%%%%%%%%%%%%%%%%%%%%%%%%%%%%%%
\section{Worked Examples}
\label{sec:feh-examples}
%%%%%%%%%%%%%%%%%%%%%%%%%%%%%%%%%%%%%%%%%%%%%%%%%%%%%%%%%%%%%%%%%%%%%%%%%%%%%%%

%------------------------------------------------------------------------------
\subsection{Example 1: Pricing Fairness in Retail}
%------------------------------------------------------------------------------

\paragraph{Setup.} A retailer considers raising prices after a supply shock.
Standard economics suggests raising prices to market-clearing level.

\paragraph{Application.} Using Fehr-Schmidt parameters:
\begin{itemize}[nosep]
    \item Customers with $\alpha > 0$ perceive price increase as ``unfair''
          (retailer gains at their expense)
    \item Loss aversion amplifies: price increase feels like a loss
    \item Expected backlash: reduced purchases, negative word-of-mouth
\end{itemize}

\paragraph{Result.} \citet{kahneman1986fairness} (with Fehr's framework)
show that ``unfair'' price increases reduce demand beyond standard elasticity.

\paragraph{Practical Implication.} Frame price increases as cost pass-through,
not profit-taking. Timing matters: increase prices when competitors do.

%------------------------------------------------------------------------------
\subsection{Example 2: Public Goods with Punishment}
%------------------------------------------------------------------------------

\paragraph{Setup.} A team faces a free-rider problem. Without intervention,
cooperation decays to 20-40\%.

\paragraph{Application.} Introduce costly peer punishment option (Fehr's
public goods game with punishment):
\begin{itemize}[nosep]
    \item Strong reciprocators (10-30\%) punish free-riders
    \item Free-riders increase contributions to avoid punishment
    \item Cooperation stabilizes at 80-95\%
\end{itemize}

\paragraph{Result.} The punishment mechanism leverages preference heterogeneity---a
minority of strong reciprocators disciplines the entire group.

\paragraph{Practical Implication.} Design peer feedback systems; make non-contribution
visible; enable low-cost ``punishment'' (social disapproval).


%%%%%%%%%%%%%%%%%%%%%%%%%%%%%%%%%%%%%%%%%%%%%%%%%%%%%%%%%%%%%%%%%%%%%%%%%%%%%%%
\section{Implications for Practice}
\label{sec:feh-practice}
%%%%%%%%%%%%%%%%%%%%%%%%%%%%%%%%%%%%%%%%%%%%%%%%%%%%%%%%%%%%%%%%%%%%%%%%%%%%%%%

\begin{tcolorbox}[colback=green!5!white, colframe=green!75!black,
    title=Practical Implications from Fehr's Research]

\textbf{1. Fairness Matters for Compliance}
\begin{itemize}[nosep]
    \item People reject ``unfair'' offers even at personal cost
    \item Frame policies as fair (procedural \& distributive justice)
    \item Anticipate resistance to perceived exploitation
\end{itemize}

\textbf{2. Leverage Heterogeneity}
\begin{itemize}[nosep]
    \item Not everyone is selfish---design for the cooperative majority
    \item Identify and empower ``strong reciprocators'' as norm enforcers
    \item Segment interventions by preference type
\end{itemize}

\textbf{3. Context is Manipulable}
\begin{itemize}[nosep]
    \item Professional vs. personal identity affects behavior
    \item Institutional framing (market vs. community) shifts norms
    \item Use $\Psi$-interventions to activate cooperative mindsets
\end{itemize}

\textbf{4. Reciprocity as Intervention Tool}
\begin{itemize}[nosep]
    \item ``Gifts'' (above-minimum wages, unexpected bonuses) elicit reciprocity
    \item Sequence matters: give first, then ask
    \item Personalize gifts for stronger effect
\end{itemize}

\textbf{5. Combine Mechanisms (Complementarity)}
\begin{itemize}[nosep]
    \item Trust + contracts > either alone
    \item Repetition + competition = super-additive cooperation
    \item Bundle interventions for multiplicative effects
\end{itemize}

\end{tcolorbox}


%%%%%%%%%%%%%%%%%%%%%%%%%%%%%%%%%%%%%%%%%%%%%%%%%%%%%%%%%%%%%%%%%%%%%%%%%%%%%%%
\section{Paper Bibliography}
\label{sec:feh-papers}
%%%%%%%%%%%%%%%%%%%%%%%%%%%%%%%%%%%%%%%%%%%%%%%%%%%%%%%%%%%%%%%%%%%%%%%%%%%%%%%

\subsection{Integration Statistics}

\begin{table}[htbp]
\centering
\caption{Fehr Paper Integration Summary}
\label{tab:feh-integration}
\begin{tabular}{@{}lr@{}}
\toprule
\textbf{Category} & \textbf{Count} \\
\midrule
Total publications & 343 \\
Integrated in EBF & 144 \\
Not integrated & 199 \\
Integration rate & 42\% \\
\midrule
\multicolumn{2}{@{}l}{\textit{Non-integrated by category:}} \\
\quad Neuroeconomics & 42 \\
\quad Pure Methodology & 18 \\
\quad German Language & 31 \\
\quad Working Papers (superseded) & 47 \\
\quad Narrow Domain & 36 \\
\quad Book Chapters/Reviews & 25 \\
\bottomrule
\end{tabular}
\end{table}

\subsection{Key Integrated Papers}

\begin{enumerate}
    \item \textbf{Fehr \& Schmidt (1999)} ``A Theory of Fairness, Competition,
          and Cooperation'' \textit{QJE}. \textbf{Citations:} 19,000+.
          \textbf{EBF use:} Foundation for $\alpha$, $\beta$ parameters.

    \item \textbf{Fehr, Fischbacher, \& Gächter (2002)} ``Strong Reciprocity,
          Human Cooperation, and the Enforcement of Social Norms''
          \textit{Human Nature}. \textbf{EBF use:} Strong reciprocity theory.

    \item \textbf{Fehr \& Gächter (1998)} ``Reciprocity and Economics''
          \textit{EER}. \textbf{EBF use:} Gift exchange parameters.

    \item \textbf{Efferson, Fehr et al. (2024)} ``Super-additive Cooperation''
          \textit{Nature}. \textbf{EBF use:} Complementarity evidence.

    \item \textbf{Fehr \& Charness (2025)} ``Social Preferences: Fundamental
          Characteristics'' \textit{JEL}. \textbf{EBF use:} Updated type distribution.
\end{enumerate}

\subsection{Non-Integrated Papers: Rationale}

Papers not integrated into EBF fall into defined categories:

\begin{tcolorbox}[colback=gray!5!white, colframe=gray!75!black,
    title=Neuroeconomics (42 papers) --- Outside EBF Scope]
\textbf{Example:} ``The Neural Basis of Altruistic Punishment'' (Science 2004)

\textbf{Rationale:} fMRI studies provide neural correlates but not extractable
behavioral parameters. Neuroeconomics is a separate domain from behavioral
economics as implemented in EBF.

\textbf{Reconsider if:} Neuroeconomics domain added to EBF.
\end{tcolorbox}


%%%%%%%%%%%%%%%%%%%%%%%%%%%%%%%%%%%%%%%%%%%%%%%%%%%%%%%%%%%%%%%%%%%%%%%%%%%%%%%
\section{Summary}
\label{sec:feh-summary}
%%%%%%%%%%%%%%%%%%%%%%%%%%%%%%%%%%%%%%%%%%%%%%%%%%%%%%%%%%%%%%%%%%%%%%%%%%%%%%%

\begin{tcolorbox}[colback=green!10!white, colframe=green!75!black,
    title=\textbf{Appendix FEH: Summary}]

\textbf{Core Question:} How do Ernst Fehr's contributions enable behavioral
economics in EBF?

\textbf{Main Contribution:}
\begin{itemize}[nosep]
    \item Established social preferences as measurable utility components
    \item Provided calibrated parameters ($\alpha=0.85$, $\beta=0.315$)
    \item Documented population heterogeneity in preferences
    \item Demonstrated complementarity in cooperation mechanisms
\end{itemize}

\textbf{Key Outputs:}
\begin{itemize}[nosep]
    \item 6 extractable parameters with confidence intervals
    \item 4-type population distribution for segmentation
    \item 2 worked examples for practitioner application
    \item 144 integrated papers with evaluation rationale
\end{itemize}

\textbf{Integration:} This appendix connects to C (CORE-WHAT), B (CORE-HOW),
V (CORE-WHEN), AAA (CORE-WHO), and BBB (CORE-WHERE) via parameter calibration
and theoretical foundations.

\end{tcolorbox}


%%%%%%%%%%%%%%%%%%%%%%%%%%%%%%%%%%%%%%%%%%%%%%%%%%%%%%%%%%%%%%%%%%%%%%%%%%%%%%%
\section{Glossary of Symbols}
\label{sec:feh-glossary}
%%%%%%%%%%%%%%%%%%%%%%%%%%%%%%%%%%%%%%%%%%%%%%%%%%%%%%%%%%%%%%%%%%%%%%%%%%%%%%%

\begin{tcolorbox}[colback=blue!5!white, colframe=blue!75!black]
\textbf{Central Reference:} For complete symbol definitions, see
\textbf{Appendix GLS} (Glossary and Notation).
\end{tcolorbox}

\begin{table}[htbp]
\centering
\caption{Symbols in Appendix FEH}
\label{tab:feh-symbols}
\renewcommand{\arraystretch}{1.3}
\begin{tabular}{@{}clp{6cm}c@{}}
\toprule
\textbf{Symbol} & \textbf{Name} & \textbf{Definition} & \textbf{In GLS?} \\
\midrule
$\alpha$ & Disadvantageous IA & Sensitivity to being behind others & \checkmark \\
$\beta$ & Advantageous IA & Sensitivity to being ahead of others & \checkmark \\
$U^S$ & Social Utility & Utility from others' outcomes & \checkmark \\
$\gamma$ & Complementarity & Interaction effect between mechanisms & \checkmark \\
$\Psi$ & Context vector & Context dimensions affecting behavior & \checkmark \\
\bottomrule
\end{tabular}
\end{table}


%%%%%%%%%%%%%%%%%%%%%%%%%%%%%%%%%%%%%%%%%%%%%%%%%%%%%%%%%%%%%%%%%%%%%%%%%%%%%%%
\section{Critical Foundations}
\label{sec:feh-foundations}
%%%%%%%%%%%%%%%%%%%%%%%%%%%%%%%%%%%%%%%%%%%%%%%%%%%%%%%%%%%%%%%%%%%%%%%%%%%%%%%

The following foundational elements underpin Fehr's contributions to EBF:

\begin{enumerate}
    \item \textbf{Methodological Foundation:} The use of incentive-compatible
    laboratory experiments with real monetary stakes, establishing behavioral
    economics as an empirical science (1990s experimental revolution).

    \item \textbf{Theoretical Foundation:} The Fehr-Schmidt utility function
    extends classical expected utility by adding social comparison terms,
    maintaining tractability while explaining cooperation puzzles.

    \item \textbf{Empirical Foundation:} Meta-analyses of 100+ ultimatum game
    studies provide stable parameter estimates ($\alpha$, $\beta$) across cultures.

    \item \textbf{Policy Foundation:} Gift exchange findings inform efficiency
    wage theory and fair wage policies in organizational economics.

    \item \textbf{EBF Integration:} Parameters from Fehr's work populate
    Appendix BBB (CORE-WHERE) and calibrate models in Appendix C (CORE-WHAT).
\end{enumerate}


%%%%%%%%%%%%%%%%%%%%%%%%%%%%%%%%%%%%%%%%%%%%%%%%%%%%%%%%%%%%%%%%%%%%%%%%%%%%%%%
\section{Open Issues}
\label{sec:feh-open-issues}
%%%%%%%%%%%%%%%%%%%%%%%%%%%%%%%%%%%%%%%%%%%%%%%%%%%%%%%%%%%%%%%%%%%%%%%%%%%%%%%

The following issues remain unresolved and represent areas for future research:

\begin{itemize}
    \item \textbf{OI-1 (Cross-cultural Stability):} While $\alpha$ and $\beta$
    estimates are robust in WEIRD populations, systematic variation across
    non-WEIRD cultures requires further validation (see Appendix BJ).

    \item \textbf{OI-2 (Dynamic Preferences):} How do social preferences evolve
    over the life course? Integration with Sutter's developmental research needed.

    \item \textbf{OI-3 (Neural Mechanisms):} Neuroeconomics papers (not integrated)
    suggest biological substrates, but EBF lacks formal incorporation of neural data.

    \item \textbf{OI-4 (Online Behavior):} Most experiments predate digital
    platforms; social preferences in online contexts remain under-explored.

    \item \textbf{OI-5 (AI Interactions):} How do humans apply fairness norms
    when interacting with AI agents? Emerging research area with EBF implications.
\end{itemize}


%%%%%%%%%%%%%%%%%%%%%%%%%%%%%%%%%%%%%%%%%%%%%%%%%%%%%%%%%%%%%%%%%%%%%%%%%%%%%%%
\section{References}
\label{sec:feh-references}
%%%%%%%%%%%%%%%%%%%%%%%%%%%%%%%%%%%%%%%%%%%%%%%%%%%%%%%%%%%%%%%%%%%%%%%%%%%%%%%

\begin{tcolorbox}[colback=blue!5!white, colframe=blue!75!black]
\textbf{Master Bibliography:} For complete reference details, see
\texttt{bcm\_master.bib} in the repository.
\end{tcolorbox}

\subsection{Key References for This Appendix}

\begin{itemize}[nosep]
    \item \textbf{Inequity Aversion:} \citet{fehr1999theory}
    \item \textbf{Strong Reciprocity:} \citet{fehr2002strong, fehr2002altruistic}
    \item \textbf{Gift Exchange:} \citet{fehr1993fairness, fehr1998reciprocity}
    \item \textbf{Context Effects:} \citet{cohn2014business, cohn2015evidence}
    \item \textbf{Recent Synthesis:} \citet{fehr2025social}
\end{itemize}

\subsection{Data Sources}

\begin{itemize}[nosep]
    \item \textbf{Google Scholar Profile:} \url{https://scholar.google.com/citations?user=WoSILroAAAAJ}
    \item \textbf{University of Zurich:} \url{https://www.econ.uzh.ch/en/people/faculty/fehr.html}
    \item \textbf{Research.com Ranking:} \url{https://research.com/u/ernst-fehr}
\end{itemize}


%%%%%%%%%%%%%%%%%%%%%%%%%%%%%%%%%%%%%%%%%%%%%%%%%%%%%%%%%%%%%%%%%%%%%%%%%%%%%%%
% CROSS-REFERENCE FOOTER
%%%%%%%%%%%%%%%%%%%%%%%%%%%%%%%%%%%%%%%%%%%%%%%%%%%%%%%%%%%%%%%%%%%%%%%%%%%%%%%

\vspace{1em}
\hrule
\vspace{0.5em}

\noindent\textit{Cross-references:}
Appendix GLS (Central Glossary),
Appendix C (CORE-WHAT),
Appendix B (CORE-HOW),
Appendix V (CORE-WHEN),
Appendix BBB (CORE-WHERE),
Appendix FRM (LIT-FEHR-METHOD).

\noindent\textit{Registry Reference:} \texttt{data/researcher-registry.yaml} (RES-FEHR-E)

\noindent\textit{Master Bibliography:} \texttt{bcm\_master.bib}

% =============================================================================
% END OF APPENDIX FEH
% =============================================================================
